\clearpage\corrigesChapitre{17}{Réflexion et réfraction}
\corrigesSouples
\begin{blocLivre}{cor-17-1}\begin{corrige}{1}{La bonne référence angulaire}
\begin{questions}
\item $i=35^\circ$.
\item $90-35=55^\circ$. La surface et la normale sont perpendiculaires.
\end{questions}
\end{corrige}\end{blocLivre}
\begin{blocLivre}{cor-17-2}\begin{corrige}{2}{Réfléchir un rayon}
\schema[corrige]{reflexion}
\begin{questions}
\item $i'=40^\circ$ par la loi de réflexion.
\item Il est de l’autre côté de la normale que le rayon incident, dans le plan d’incidence.
\end{questions}
\end{corrige}\end{blocLivre}
\begin{blocLivre}{cor-17-3}\begin{corrige}{3}{Lire un indice}
\begin{questions}
\item Non, l’indice est sans unité.
\item Oui pour une incidence non nulle : le milieu d’arrivée a l’indice le plus élevé.
\end{questions}
\end{corrige}\end{blocLivre}
\begin{blocLivre}{cor-17-4}\begin{corrige}{4}{Une incidence normale}
\begin{questions}
\item $i=0^\circ$ car le rayon suit la normale.
\item $\sin r=0$, donc $r=0^\circ$ : le rayon transmis n’est pas dévié.
\end{questions}
\end{corrige}\end{blocLivre}
\begin{blocLivre}{cor-17-5}\begin{corrige}{5}{Calculer un angle réfracté}
\schema[corrige]{refraction}
\begin{questions}
\item $\sin r=(1,00/1,50)\sin30,0^\circ=0,3333\ldots$.
\item $r\simeq19,5^\circ<i$ : le rayon se rapproche de la normale.
\end{questions}
\end{corrige}\end{blocLivre}
\begin{blocLivre}{cor-17-6}\begin{corrige}{6}{Retrouver un indice}
\begin{questions}
\item $n=\sin i/\sin r$ avec l’indice de l’air pris égal à 1.
\item $n=\sin45,0^\circ/\sin28,1^\circ\simeq1,50$.
\end{questions}
\end{corrige}\end{blocLivre}
\begin{blocLivre}{cor-17-7}\begin{corrige}{7}{Lire une droite de sinus}
\schema[corrige]{snell-graphe}
\begin{questions}
\item Par exemple, le point $(0,40;0,60)$ donne une pente de $0,60/0,40=1,50$.
\item $n\simeq1,50$ car la pente vaut $n/n_{\mathrm{air}}$.
\end{questions}
\end{corrige}\end{blocLivre}
\begin{blocLivre}{cor-17-8}\begin{corrige}{8}{Pourquoi un prisme disperse}
\schema[corrige]{prisme}
\begin{questions}
\item Les lois de réfraction donnent des déviations différentes lorsque les indices diffèrent.
\item Non : il sépare les radiations présentes dans la lumière incidente.
\end{questions}
\end{corrige}\end{blocLivre}
\begin{blocLivre}{cor-17-9}\begin{corrige}{9}{Une sortie vers l’air}
\begin{questions}
\item $\sin r=1,50\sin20,0^\circ\simeq0,513$ ; $r\simeq30,9^\circ$.
\item Oui : $r>i$, cohérent avec le passage vers l’indice plus faible.
\end{questions}
\end{corrige}\end{blocLivre}
\begin{blocLivre}{cor-17-10}\begin{corrige}{10}{Une pente inversée}
\begin{questions}
\item Les axes sont inversés par rapport au graphique précédent : la pente vaut $n_{\mathrm{air}}/n_{\mathrm{verre}}$.
\item $n_{\mathrm{verre}}\simeq1/0,667\simeq1,50$.
\end{questions}
\end{corrige}\end{blocLivre}
\begin{blocLivre}{cor-17-11}\begin{corrige}{11}{Un angle impossible}
\begin{questions}
\item Le sinus d’un angle réel est compris entre -1 et 1.
\item Aucun rayon réfracté propagatif ne satisfait cette équation : on rencontre la réflexion totale. Cette situation est un prolongement guidé, pas une nouvelle formule à mémoriser.
\end{questions}
\end{corrige}\end{blocLivre}
\begin{blocLivre}{cor-17-12}\begin{corrige}{12}{Tester une série}
\demarche{Comparer plusieurs rapports et les écarts à la précision des mesures.}
\begin{questions}
\item Les rapports valent 1,50 ; 1,50 ; environ 1,53. Ils sont proches.
\item Non : il faut comparer les écarts à la précision des mesures et du tracé, et rechercher des erreurs systématiques.
\end{questions}
\retiens{Un modèle expérimental se juge à l’échelle des incertitudes, pas sur l’égalité parfaite de nombres arrondis.}
\end{corrige}\end{blocLivre}

